\begin{proposition}[7.4.6]
\label{IV.2.s7.4.6-fr}
Supongamos que la propiedad $\boldsymbol P$ satisface las condiciones
$(\mathrm P_{\mathrm I})$, $(\mathrm P_{\mathrm{II}})$ y
$(\mathrm P_{\mathrm{III}})$. Sean $A$ un anillo noetheriano,
$\mathfrak J$ un ideal de $A$ y $\widehat A$ el completado separado de
$A$ para la topología $\mathfrak J$-preádica. Entonces, si $A$ es un
$\boldsymbol P$-anillo (7.4.5), el morfismo canónico
$\operatorname{Spec}(\widehat A)\to\operatorname{Spec}(A)$ es un
$\boldsymbol P$-morfismo.
\end{proposition}

\begin{proof}
Utilizando (7.3.2, \emph{c$'$)}), basta demostrar que, para todo ideal
maximal $\mathfrak n$ de $\widehat A$ cuya imagen inversa es
$\mathfrak m$ en $A$, el morfismo
$\operatorname{Spec}((\widehat A)_{\mathfrak n})\to
\operatorname{Spec}(A_{\mathfrak m})$ es un $\boldsymbol P$-morfismo.
Se sabe (Bourbaki, \emph{Alg. comm.}, cap. III, \S~3,
n\textsuperscript{o}~4, prop. 8) que el homomorfismo canónico
$A_{\mathfrak m}\to(\widehat A)_{\mathfrak n}$ es inyectivo, que la
topología $\mathfrak mA_{\mathfrak m}$-preádica sobre $A_{\mathfrak m}$
está inducida por la topología
$\mathfrak n(\widehat A)_{\mathfrak n}$-preádica y que
$A_{\mathfrak m}$ es denso en $(\widehat A)_{\mathfrak n}$ para esta
topología; por consiguiente, el completado de $A_{\mathfrak m}$ para la
topología $\mathfrak mA_{\mathfrak m}$-preádica se identifica con el de
$(\widehat A)_{\mathfrak n}$ para la topología
$\mathfrak n(\widehat A)_{\mathfrak n}$-preádica. Se tienen, pues, dos
morfismos
\[
\operatorname{Spec}((A_{\mathfrak m})^{\widehat{\,}})
\xrightarrow{f}\operatorname{Spec}((\widehat A)_{\mathfrak n})
\xrightarrow{g}\operatorname{Spec}(A_{\mathfrak m})
\]

\oldpage[IV-2]{203}

tales que $f$ es fielmente plano; como, por hipótesis, $g\circ f$ es un
$\boldsymbol P$-morfismo, también lo es $g$ en virtud de
$(\mathrm P_{\mathrm{II}})$.
\end{proof}
